\documentclass{article}
\usepackage{main}

\begin{document}



\begin{enumerate}[left=0pt,itemsep=10pt]

\item $105$ minutes correspondent à :\\



\begin{tabular}{l@{\qquad}l@{\qquad}l@{\qquad}l}

a. $1,05h$ & b. $1,45h$ & c. $1,5h$ & d. $1,75h$ \\

\end{tabular}



\vspace{0.8cm}



\item On considère un nombre réel $x$. Si $x>0$, alors : \\



\begin{tabular}{l@{\qquad}l@{\qquad}l@{\qquad}l}

a. $x=x^2$ & b. $x \leq x^2$ & c. $x \geq x^2$ & d. Cela dépend de la valeur de $x$ \\

\end{tabular}



\vspace{0.8cm}



\item On considère la relation : $A=\dfrac{x+6y}{x}$. Lorsque $x=-\dfrac{1}{3}$ et $y=\dfrac{1}{18}$, la valeur de $A$ est : \\



\begin{tabular}{l@{\qquad}l@{\qquad}l@{\qquad}l}

a. $\dfrac{1}{3}$ & b. $0$ & c. $3$ & d. $-{\dfrac{1}{3}}$

\end{tabular}



\vspace{0.8cm}



\item Le volume d'un cylindre de hauteur $h$ et de rayon $r$ est donné par $V=\pi r^2 h$. On cherche à isoler $h$. On a : \\



\begin{tabular}{l@{\qquad}l@{\qquad}l@{\qquad}l}

a. $h=\sqrt{\dfrac{V}{\pi r^2}}$ 

& b. $h=\dfrac{\pi r^2}{V}$ 

& c. $h=\dfrac{V}{\pi r^2}$ 

& d. $h=\dfrac{r^2}{\pi V}$

\end{tabular}



\vspace{0.8cm}



\item Un drapeau a une aire de $3\,m^2$. Cela représente :



\begin{tabular}{l@{\qquad}l@{\qquad}l@{\qquad}l}

a. $0{,}003\,cm^2$ 

& b. $0{,}3\,cm^2$ 

& c. $300\,cm^2$ 

& d. $30\,000\,cm^2$ \\

\end{tabular}



\vspace{0,8cm}



\item Le prix d'un article est multiplié par $1,11$. Cela signifie que le prix de l'article a subi : 



\begin{tabular}{l@{\qquad}l@{\qquad}l@{\qquad}l}

a. une hausse de $11$ euros

& b. une hausse de $11\%$ \\

c. une hausse de $1{,}1\%$ 

& d. une hausse de $111\%$ \\

\end{tabular}



\vspace{0,8cm}



\item  Une forme factorisée de l'expression $-3(x+1)+(x+1)^2$ est :



\begin{tabular}{l@{\qquad}l@{\qquad}l@{\qquad}l}

a. $-(x+1)$ & b. $-(x+1)(x+2)$ & c. $(x+1)(x-2)$ & d. $x^2-x-2$ \\

\end{tabular}





\vspace{0,8cm}



\item Le point qui n'appartient pas à la courbe d'équation $y=\dfrac{6}{x}$ est : \\

\begin{tabular}{l@{\qquad}l@{\qquad}l@{\qquad}l}

a. $M(2;3)$ & b. $N(3;3)$ & c. $P(1;6)$ & d. $Q(0,5;12)$ \\

\end{tabular}



\vspace{0,8cm}



\item Une réduction de $50\%$ suivie d'une augmentation de $50\%$ équivaut à : \\

\begin{tabular}{l@{\qquad}l@{\qquad}l@{\qquad}l}

a. revenir à la valeur initiale & b. une baisse de $25\%$ & \\ c. une baisse de $75\%$ & d. une hausse de $75\%$ \\

\end{tabular}



\vspace{0,8cm}



\item Une boite contient 20 billes au total. Dans cette boite, la proportion de billes rouges est égale à 0,3. Le nombre de billes rouges dans la boite est égal à : \\

\begin{tabular}{l@{\qquad}l@{\qquad}l@{\qquad}l}

a. $6$ & b. $5$ &  c. $3$ & d. $17$ \\

\end{tabular}



\vspace{0,8cm}



\item  On peut affirmer que : \\

\begin{tabular}{l@{\qquad}l@{\qquad}l@{\qquad}l}

a. $\dfrac{12}{13}<\dfrac{13}{14}$ & b. $\dfrac{12}{13}=\dfrac{13}{14}$ & c. $\dfrac{12}{13}>\dfrac{13}{14}$ & d. Il est impossible de comparer ces deux nombres. \\

\end{tabular}



\vspace{0,8cm}



\item  La moitié de $10^{24}$ vaut :  \\

\begin{tabular}{l@{\qquad}l@{\qquad}l@{\qquad}l}

a. $5^{24}$ & b. $10^{12}$ & c. $5^{12}$ & d. $5 \times 10^{23}$ \\

\end{tabular}



\end{enumerate}



\end{document}